\documentclass[10pt,a4paper]{amsart}
\usepackage[margin=19mm]{geometry}
\usepackage{amsmath,amssymb,mathtools}
\usepackage[scr=rsfs,cal=euler]{mathalfa}
\usepackage{xfrac}
\usepackage[unicode,hidelinks,bookmarks=false]{hyperref}
\newcommand{\ie}{i.e.\ }
\newcommand{\cf}{cf.\ }
\newcommand{\resp}{respectively\ }
\newcommand{\Subsection}{\subsection}
\providecommand{\nobreakdash}{\nobreak\hbox{-}}
\setlength{\emergencystretch}{2em}
\multlinegap0pt
\newtheorem{example}{Example}
\newtheorem{lemma}{Lemma}
\newcommand{\ba}{\begin{array}}
\newcommand{\bb}{\begin{equation}}
\newcommand{\ea}{\end{array}}
\newcommand{\ee}{\end{equation}}
\title{A look on equations describing pseudospherical surfaces}
\author{Igor Leite Freire}
\date{Source: arXiv:2506.23890v4 (CC BY 4.0)}
\begin{document}
\maketitle

\noindent\emph{Source and licence.} A look on equations describing pseudospherical surfaces. Version arXiv:2506.23890v4 (CC BY 4.0), distributed by arXiv under the Creative Commons Attribution 4.0 International licence, http://creativecommons.org/licenses/by/4.0/, as recorded in the arXiv licence metadata for this exact version and shown on the abstract page https://arxiv.org/abs/2506.23890v4. Full licence notice: \texttt{licenses/arxiv-2506.23890-CC-BY-4.0.txt}.\par\smallskip

\noindent\textit{Editorial scope.} Counted unit: the article's lemma lema3.1 (source0.tex lines 564--566). The article's complete original proof of it is delivered with it (source0.tex lines 568--570, one \texttt{proof} environment of 3 lines). No mathematics is written, repaired or re-derived: the statement and the proof are the article's own lines, in the article's own notation, and no retained line is substituted or re-flowed.  Retained uncounted as the source-local context the proof needs: lines 316--370.  Nothing was omitted from any retained span, so the package carries no omission record.  No retained line contains a citation, so the wrapper carries no bibliography and no bibliography entry.  No retained line contains a figure environment, an image-inclusion command or drawing code, so no image is delivered and the package declares no asset dependency.  Editorial formatting only, in addition: an AMS \texttt{amsart} preamble that restates the article's own declarations and macros used by the retained lines, namely the environments \texttt{example}, \texttt{lemma} on separate counters, together with the standard packages those lines need.  The retained environments print in this document as Example 1, Lemma 1 in order of appearance, whereas the article prints the same items with the numbering of its own full document.  The complete native source of the article as distributed in the arXiv e-print for this exact version is bundled unchanged as \texttt{sources/180745/source0.tex}.
\begin{example}\label{ex2.2}
     Let $m=u-u_{xx}$, $\lambda\in\mathbb{R}\setminus\{0\}$ and consider the one-forms
\bb\label{2.0.8}
\ba{lcl}
\omega_1&=&{\Big(\frac{\lambda}{2}+\frac{1}{2\lambda}-m\Big)dx+\Big(um+\frac{\lambda}{2}u-\frac{u}{2\lambda}-\frac{1}{2}-\frac{\lambda^2}{2}\Big)dt},\\
\\
\omega_2&=&-u_xdt,\\
\\
\omega_3&=&{\Big(m+\frac{1}{2\lambda}-\frac{\lambda}{2}\Big)dx+\Big(\frac{\lambda^2}{2}-\frac{1}{2}-\frac{u}{2\lambda}-\frac{\lambda}{2}u-um\Big)dt}.
\ea
\ee
A straightforward calculation shows that
\bb\label{2.0.9}
\ba{lcl}
d\omega_1-\omega_3\wedge\omega_2&=&\Big(u_t-u_{txx}+3uu_x-2u_xu_{xx}-uu_{xxx}\Big)dx\wedge dt,\\
\\
d\omega_2-\omega_1\wedge\omega_3&=&0,\\
\\
d\omega_3-\omega_1\wedge\omega_2&=&-\Big(u_t-u_{txx}+3uu_x-2u_xu_{xx}-uu_{xxx}\Big)dx\wedge dt,
\ea
\ee
meaning that on the generic solutions of the CH equation 
   \bb\label{1.0.5}
    u_t-u_{txx}+3uu_x=2u_xu_{xx}+uu_{xxx},
    \ee
    or its equivalent form
    \bb\label{1.0.6}
    m_t+2um_x+u_xm=0,
    \ee
the one-forms \eqref{2.0.8} define a two-dimensional Riemannian manifold with first fundamental form
$$
\ba{lcl}
I&=&{\Big[m^2-\Big(\lambda+\frac{1}{\lambda}\Big)m+\frac{1}{4}\Big(\lambda+\frac{1}{\lambda}\Big)^2\Big]dx^2}+\\
\\
&&{\Big[-um^2+\frac{1}{\lambda}um+\frac{1}{4}\Big(\lambda^2-\frac{1}{\lambda^2}\Big)u+\lambda\Big(\frac{\lambda}{2}+\frac{1}{2\lambda}\Big)m-\lambda\Big(\frac{\lambda}{2}+\frac{1}{2\lambda}\Big)^2\Big]dxdt}+\\
\\
&&{\Big[um^2+u^2m\Big(\lambda-\frac{1}{\lambda}\Big)+u_x^2-\lambda\Big(\lambda+\frac{1}{\lambda}\Big)um+\frac{1}{4}\Big(\lambda+\frac{1}{\lambda}\Big)^2u^2-\frac{\lambda}{2}\Big(\lambda+\frac{1}{\lambda}\Big)^2}\\
\\
&&{+\frac{\lambda^2}{4}\Big(\lambda+\frac{1}{\lambda}\Big)^2\Big]dt^2}.
\ea
$$

Let us investigate the generic solutions of the CH equation. It is straightforward to see that
   $$
    \omega_1\wedge\omega_2=-\Big(\frac{\lambda}{2}+\frac{1}{2\lambda}-m\Big)u_x dx\wedge dt.
   $$
    Provided that $u$ is a solution of the CH equation, then $\omega_1\wedge\omega_2=0$ somewhere if and only if 
    \bb\label{2.0.10}
    m=\frac{\lambda}{2}+\frac{1}{2\lambda}
    \ee
    or
    \bb\label{2.0.11}
    u_x=0.
    \ee
\end{example}

\begin{lemma}\label{lema3.1}
    Assume that $u\in C^{0}(H^{4}(\mathbb{R}),[0,T))\cap C^{1}(H^{3}(\mathbb{R}),[0,T))$ is a non-trivial solution of the CH equation \eqref{1.0.6}. If $u_x\neq 0$ on some open set $\Omega$, then we cannot have $m=c\neq 0$, where $c$ is a constant.
\end{lemma}

\begin{proof}
    Suppose $m=c\neq 0$ is a constant on some open and non-empty set $\Omega$, then necessarily $m\neq 0$ and \eqref{1.0.6} implies \eqref{2.0.11}. On the other hand, if $u$ is a solution of \eqref{1.0.6} such that $u_x=0$ on an open set $\Omega$, then $u_x=u_{xx}=u_{xxx}=u_{txx}=0$ and $u_t=0$ on $\Omega$, which, once substituted into \eqref{1.0.6}, implies $u(x,t)=c$.
\end{proof}

\end{document}
